% ===== Appendix: corrections to the book, what we added, sources (after Questions; not presented) =====
\appendix
\partaccent{lavender}\gdef\idiottopicemoji{books}\gdef\idiottopicname{Հավելված}
\newcolumntype{L}[1]{>{\raggedright\arraybackslash}p{#1}}

\begin{frame}{Հավելված Ա. Ուղղումներ գրքի նկատմամբ (1/2)}
  \scriptsize
  \setlength{\tabcolsep}{4pt}
  \renewcommand{\arraystretch}{1.3}
  \begin{tabular}{@{}L{2.3cm}L{3.3cm}L{4.6cm}L{3.3cm}@{}}
    \toprule
    \textbf{Թեմա} & \textbf{Գրքում (2016)} & \textbf{Այսօր} & \textbf{Աղբյուր}\\
    \midrule
    Աշի փորձը & «75\% դեպքերում սխալ էին»; «որն է ամենաերկարը» & Սխալները մոտ 37\% էին, 75\%-ը սխալվել է գոնե մեկ անգամ։ Հարցը՝ որն է նույն երկարությամբ & Asch, Psychol. Monogr., 1956\\
    Միլգրամ & 65\%-ը գնաց մինչև վերջ & Առաջին տարբերակում՝ այո։ 21 տարբերակի միջինը 43,6\% է & Haslam et al., PLoS ONE, 2014; Perry, 2012\\
    Ստենֆորդի բանտ & Պահակներն իրենք դարձան դաժան & Նրանց հրահանգել էին։ BBC-ի նման փորձում չկրկնվեց & Le Texier, Am. Psychol., 2019; Reicher \& Haslam, 2006\\
    Ամբոխ & Կորցնում ենք անհատականությունը & Ավելի խիստ հետևում ենք խմբի նորմերին & Postmes \& Spears, Psychol. Bull., 1998\\
    Դիտորդի էֆեկտ & 38 անտարբեր վկա (գիրքն ինքն է նշում չափազանցությունը) & Իրական կոնֆլիկտների 91\%-ում գոնե մեկն օգնում է & Philpot et al., Am. Psychol., 2020\\
    \bottomrule
  \end{tabular}
\end{frame}

\begin{frame}{Հավելված Ա. Ուղղումներ գրքի նկատմամբ (2/2)}
  \scriptsize
  \setlength{\tabcolsep}{4pt}
  \renewcommand{\arraystretch}{1.3}
  \begin{tabular}{@{}L{2.3cm}L{3.3cm}L{4.6cm}L{3.3cm}@{}}
    \toprule
    \textbf{Թեմա} & \textbf{Գրքում (2016)} & \textbf{Այսօր} & \textbf{Աղբյուր}\\
    \midrule
    Դիմախաղ & Բնածին է ու համընդհանուր & Վիճելի է. ազատ դասավորման դեպքում հեռավոր մշակույթը այլ կերպ է խմբավորում & Gendron et al., Emotion, 2014; Barrett et al., 2019\\
    Օքսիտոցին & Կապված է վստահության հետ & Քթի սփրեյով փորձերը լավ չեն կրկնվում & Nave et al., Perspect. Psychol. Sci., 2015\\
    Բաժանում ու ցավ & Նույն շրջանները, ինչ ֆիզիկական ցավը & Մասամբ համընկնում են, բայց ակտիվության նախշերը տարբեր են & Woo et al., Nat. Commun., 2014\\
    Սոցիալական ուղեղ & Մեծ ուղեղը՝ սոցիալական կյանքից & Պրիմատների մոտ սննդակարգն ավելի լավ է կանխատեսում ուղեղի չափը. վեճը բաց է & DeCasien et al., Nat. Ecol. Evol., 2017\\
    Ցածր գնդակ & Նախ համաձայնել էին ժամը 9-ին & Սկզբում ժամը չէին ասում (56\% ընդդեմ 31\%) & Cialdini et al., JPSP, 1978\\
    Դուռը երեսին & «Ուսանողներ» & 6-8 տարեկան երեխաներ & Chan \& Au, 2011\\
    \bottomrule
  \end{tabular}
\end{frame}

\begin{frame}{Հավելված Բ. Ինչ ենք ավելացրել գրքից դուրս}
  \begin{columns}[T]
    \begin{column}{0.48\textwidth}
      \begin{card}[colback=mint!30!cardwhite, fontupper=\scriptsize]
        {\blackfont Թվեր ու հետազոտություններ}\par\smallskip
        \begin{itemize}\setlength{\itemsep}{0.5mm}
          \item Մենախցի սահմանումը՝ ՄԱԿ-ի Նելսոն Մանդելայի կանոններ, 2015
          \item «7-38-55»՝ Mehrabian \& Wiener, 1967
          \item Երեխաների տարիքը՝ Chan \& Au, 2011
          \item Ժամը 7-ի փորձի թվերը՝ Cialdini et al., 1978
          \item Աշի թվերը՝ Asch, 1956
          \item Դիտորդի էֆեկտի թվերը՝ Darley \& Latané, 1968
          \item Միլգրամի տարբերակները՝ Milgram, 1974
          \item «Գիտությունն այսօր» տուփերը (ուղղումների աղյուսակում)
        \end{itemize}
      \end{card}
    \end{column}
    \begin{column}{0.48\textwidth}
      \begin{card}[colback=peach!35!cardwhite, fontupper=\scriptsize]
        {\blackfont Մերն է, ոչ թե գրքից}\par\smallskip
        \begin{itemize}\setlength{\itemsep}{0.5mm}
          \item Աֆազիաների հայերեն օրինակները (գրքում՝ անգլերեն)
          \item «Պետք չէր»-ի, էմոտիկոնների ու Աշի չաթային «փորձերը»
          \item Գումարները դրամով (գրքում՝ ֆունտով)
          \item Լեզվի շրջանների նկարը (Wikimedia Commons, հանրային սեփականություն), սիրամարգի ու Նյուրնբերգի լուսանկարները (Wikimedia Commons)
          \item Մեմերը. բովանդակությունը գրքից է
        \end{itemize}
      \end{card}
    \end{column}
  \end{columns}
\end{frame}

\begin{frame}{Հավելված Գ. Աղբյուրներ (1/2)}
  \scriptsize
  Գրքի ծանոթագրություններից, որոնց վրա հիմնված են սլայդները\par\smallskip
  \begin{columns}[T]
    \begin{column}{0.48\textwidth}
      \begin{itemize}\setlength{\itemsep}{0.4mm}
        \item Conley A. Torture in US jails and prisons. Vienna J. Int. Const. Law, 2013
        \item Pasley B. N. et al. Reconstructing speech from human auditory cortex. PLoS Biol., 2012
        \item Lucy J. A. Language Diversity and Thought. 1992
        \item Davies I. R. A study of colour grouping in three languages. Br. J. Psychol., 1998
        \item Sacks O. The Man Who Mistook His Wife for a Hat. 1985
        \item Whalen P. J. et al. Neuroscience and facial expressions of emotion. Emotion Rev., 2013
        \item Guéguen N. Foot-in-the-door technique and computer-mediated communication. Comput. Hum. Behav., 2002
      \end{itemize}
    \end{column}
    \begin{column}{0.48\textwidth}
      \begin{itemize}\setlength{\itemsep}{0.4mm}
        \item Chan A. C., Au T. K. Getting children to do more academic work. Teach. Teach. Educ., 2011
        \item Ebster C., Neumayr B. Applying the door-in-the-face technique to retailing. 2008
        \item Burger J. M., Cornelius T. Raising the price of agreement. J. Appl. Soc. Psychol., 2003
        \item Cialdini R. B. et al. Low-ball procedure for producing compliance. JPSP, 1978
        \item Farrow T. F. et al. Neural correlates of self-deception and impression-management. Neuropsychologia, 2015
        \item Coxworth J. E. et al. Grandmothering life histories and human pair bonding. PNAS, 2015
        \item Aron A. et al. Reward, motivation, and emotion systems in early-stage romantic love. J. Neurophysiol., 2005
      \end{itemize}
    \end{column}
  \end{columns}
\end{frame}

\begin{frame}{Հավելված Դ. Աղբյուրներ (2/2)}
  \scriptsize
  Գրքի ծանոթագրություններից, որոնց վրա հիմնված են սլայդները\par\smallskip
  \begin{columns}[T]
    \begin{column}{0.48\textwidth}
      \begin{itemize}\setlength{\itemsep}{0.4mm}
        \item Hays W. S. Human pheromones: have they been demonstrated? Behav. Ecol. Sociobiol., 2003
        \item Kross E. et al. Social rejection shares somatosensory representations with physical pain. PNAS, 2011
        \item Fisher H. E. et al. Reward, addiction, and emotion regulation systems associated with rejection in love. J. Neurophysiol., 2010
        \item Dunbar R. I. The social brain hypothesis and its implications. Ann. Hum. Biol., 2009
        \item Dávid-Barrett T., Dunbar R. Processing power limits social group size. Proc. R. Soc. B, 2013
        \item Asch S. E. Studies of independence and conformity. Psychol. Monogr., 1956
        \item Latané B., Darley J. M. Bystander «apathy». Am. Sci., 1969
      \end{itemize}
    \end{column}
    \begin{column}{0.48\textwidth}
      \begin{itemize}\setlength{\itemsep}{0.4mm}
        \item Janis I. L. Groupthink. 1982
        \item Reicher S. D., Spears R., Postmes T. A social identity model of deindividuation phenomena. Eur. Rev. Soc. Psychol., 1995
        \item Milgram S. Behavioral study of obedience. J. Abnorm. Soc. Psychol., 1963
        \item Morrison S., Decety J., Molenberghs P. The neuroscience of group membership. Neuropsychologia, 2012
        \item Zimbardo P. G. Stanford Prison Experiment. 1971
        \item Silani G. et al. Right supramarginal gyrus is crucial to overcome emotional egocentricity bias. J. Neurosci., 2013
        \item Strömwall L. A. et al. Victim blame and the just world hypothesis. J. Sex. Aggress., 2013
      \end{itemize}
    \end{column}
  \end{columns}
\end{frame}
